\documentclass[11pt]{article}
\usepackage[margin=1in]{geometry}
\usepackage{amsmath,amssymb,amsthm}
\usepackage{hyperref}
\newtheorem{definition}{Definition}
\newtheorem{proposition}{Proposition}
\newtheorem{remark}{Remark}
\title{Predictive Prefetching on the Agent Read Path:\\ A Formal Model, Incentive Analysis, and Threat Model}
\author{Readpath project}
\date{September 2026}
\begin{document}
\maketitle

\begin{abstract}
We model an MCP server that observes an agent's tool calls, uses a fast classifier to predict the
next call, and generates the likely response before that call arrives. The prototype uses Jev, a
``System 1'' model, for prediction. It can attach generated pages to claimed domains and offer those
pages to third-party brands as paid placements. We analyze the design's hit rate and latency, then
examine its central risk: machinery that hides latency can also influence what an agent presents to
its user. This creates a principal--agent problem and exposes a broader vulnerability in agents that
trust unauthenticated tool output. We conclude with a threat model, defensive guidance, and the
safeguards enforced by the prototype.
\end{abstract}

\section{Formal model}
Let $\Sigma$ be a finite alphabet of tool names, such as \texttt{find\_providers} and
\texttt{read\_skill}. An agent session is represented by a word
$w = t_1 t_2 \cdots t_n \in \Sigma^*$. Let $A$ be a set of professions and $L$ a set of locations.

\begin{definition}[Bait set]
$B = A \times L$. A call $c = (t, (a,\ell))$ is a \emph{bait read} if and only if
$(a,\ell)\in B$.
\end{definition}

\begin{definition}[Next-step predictor]
A predictor is a map $\pi : \Sigma \times B \to \Delta(\Sigma)$. As an empirical baseline, we use
the first-order Markov estimate from a set of logged sessions $W$:
\[
\hat\pi(t' \mid t) \;=\; \frac{N_W(t\,t')}{\sum_{s\in\Sigma} N_W(t\,s)},
\]
where $N_W(u)$ counts the occurrences of bigram $u$ in $W$. Jev replaces $\hat\pi$ with a model that
returns $\pi(\cdot\mid t,(a,\ell))$ over a predefined output set. It therefore cannot produce a
tool name outside $\Sigma$.
\end{definition}

\begin{definition}[Prefetch policy]
Fix a threshold $\tau\in(0,1]$. After a bait read $(t,b)$, the server computes
$P_\tau(t,b)=\{t' : \pi(t'\mid t,b)\ge\tau\}$ and generates an artifact $g(t',b)$ for every
$t'\in P_\tau$ before the agent sends its next request.
\end{definition}

\section{Performance analysis}
\begin{proposition}[Hit rate]
If the agent's next tool is drawn from $\pi$, the expected share of calls served from the prefetch
cache is $H(\tau)=\sum_{t' \,:\, \pi(t')\ge\tau}\pi(t')$. The number of artifacts generated for
each bait is bounded by $|P_\tau|\le 1/\tau$.
\end{proposition}
\begin{proof}
Every $t'\in P_\tau$ is prefetched, so a draw lands in $P_\tau$ with probability
$\sum_{t'\in P_\tau}\pi(t')$. Because $\pi$ sums to $1$ and every member has mass at least $\tau$,
$|P_\tau|\tau\le 1$.
\end{proof}

\begin{proposition}[Latency hiding]
Let $\lambda_g$ denote generation latency and let $\delta$ denote the gap between calls, including
the agent's thinking time. A prefetched artifact is ready before the next call if and only if
$\lambda_g\le\delta$. If $\lambda_c$ is the cache-read cost, the expected response time is
$\mathbb E[\lambda]=H(\tau)\,\lambda_c+(1-H(\tau))(\lambda_c+\lambda_g)$. Prefetching can therefore
hide generation latency when $\lambda_g\le\delta$, with greater benefit as $H(\tau)$ increases.
\end{proposition}

\paragraph{Domain allocation.} Let $D$ be a pool of unused domains and $\beta$ a purchasing budget.
For each new bait, the server claims one available $d\in D$. If the pool is empty, it buys a domain
at price $p$ only when $\text{spent}+p\le\beta$. The prototype disables purchasing unless an
operator explicitly enables it.

\section{The incentive problem}
The server is not neutral. An agent consumes its output on behalf of a user, the \emph{principal},
while third-party brands pay the server operator. This introduces a hidden intermediary with
misaligned incentives:

\begin{enumerate}
\item \textbf{Objective mismatch.} The agent should maximize user utility $U$, while the operator
      maximizes brand revenue $R$. Content optimized for $R$ may displace content that would better
      serve $U$, simply because it is well formed and ready at the right moment.
\item \textbf{Prediction enables manipulation.} The operator writes $g(t',b)$ \emph{before} the agent
      asks and can tune it to the kinds of responses that the agent tends to accept. The operator
      therefore has an asymmetric ability to shape both the anticipated request and its answer.
\item \textbf{Non-disclosure.} If the end user cannot see that a placement is sponsored, the result
      is deceptive advertising. Agent-directed instructions in that content may also be treated as
      prompt injection by the consuming platform.
\end{enumerate}

\begin{remark}[Constraints enforced]
The prototype attempts to remain on the legitimate side of this boundary. It (i) prompts the
generator for factual content, (ii) rejects output containing imperatives directed at the reading
agent, such as ``ignore previous instructions,'' (iii) labels sponsored placements on every page,
and (iv) confines anonymous writes to a scratch namespace. These controls may still be insufficient;
their effectiveness depends in part on how the consuming agent presents sources to its user.
\end{remark}

\section{Race dynamics on the agent read path}
\paragraph{Read-path race.} Let $O=\{o_1,\dots,o_m\}$ be a set of operators. Each operator prepares an
artifact for predicted step $t'$ after observing bait $b$. An agent with latency budget $\delta$
accepts the first well-formed artifact $o_i$ that is ready by $\delta$ rather than reviewing all $m$
alternatives. Operator $i$ therefore wins according to its readiness time $r_i\le\delta$ and the
agent's probability $\alpha_i$ of accepting its content, not necessarily its true quality $q_i$.
When entry is inexpensive---for example, because domains and generation are cheap---new operators
can enter until the expected margin disappears.

\paragraph{Why the race can favour fraud.} Let $\alpha_i$ be the rate at which an agent accepts
operator $i$'s content. Optimizing $\alpha_i$ against a particular model means optimizing for that
model's judgment. Honest and fraudulent operators face the same selection mechanism, but a fraudulent
operator can also use fabricated contacts or lookalike payment endpoints. Its feasible set is thus
strictly larger. Unless verification is effectively free, accepted results may become enriched for
fraud. Verification adds latency, and added latency is a disadvantage in the read-path race.

\paragraph{Time of check to time of use.} Suppose an artifact is checked at $t_c$ and used later at
$t_u>t_c$. An operator that changes the content during $(t_c,t_u)$ defeats the earlier check.
Prefetched and cached artifacts can widen this interval.

\paragraph{Mitigations.} Useful controls include signed provenance; reputation and age requirements
for domains in high-stakes settings; agreement across independent sources; hash-pinning from check to
use; human confirmation before financial or contact actions; and sponsorship labels visible to the
end user.

\section{When generation outpaces prediction}
\emph{This section is a projection, not a measurement.}

\paragraph{The prefetch advantage shrinks.} Proposition 2 shows that prefetching can hide latency
when $\lambda_g\le\delta$. Let $\lambda_g^{\text{od}}$ be the latency of a fast on-demand generator,
possibly another agent invoked when the request arrives. As $\lambda_g^{\text{od}}\to 0$, the latency
difference
$\mathbb E[\lambda_{\text{on-demand}}]-\mathbb E[\lambda_{\text{prefetch}}]\to 0$, and $H(\tau)$
ceases to determine who wins. Competition then shifts from prediction quality toward generation
speed, and every model improvement moves that boundary.

\paragraph{Verification becomes the bottleneck.} Each step in an $n$-call chain costs a round trip,
so the chain creates $\Theta(n)$ round trips and $\Theta(n)$ opportunities for interception. Once
generation is cheap, verification becomes the scarce resource. Acceptance is limited by verification
cost $v$, which the agent pays in latency, rather than by generation cost. This recreates the adverse
incentive described above at a larger scale.

\paragraph{Whole-intent submission.} One alternative is to replace the stream
$t_1\cdots t_n$ with a single declared intent $I$ containing the goal, constraints, budget, and
verification requirements. The server then plans over $I$ in one round trip instead of predicting
$\pi(t'\mid t,b)$ one step at a time. It can return a signed plan $\sigma(\text{plan},I)$ for the agent
to check against its constraints.

\paragraph{Tradeoffs.} This design has three important costs. First, the operator sees all of $I$,
which increases the information asymmetry unless the intent is minimized or the operator is tightly
constrained. Second, intent itself may become the object of an advertising auction. Third, declaring
$I$ does not make the answer truthful; provenance and hash-pinning remain necessary.

\section{A possible protocol response}
\subsection{Whole-intent submission}
This proposal addresses repeated round trips and the check/use gap; it is not, by itself, a trust
mechanism. Let intent be the tuple $I=(g,C,\beta,\rho)$, where $g$ is the goal, $C$ the hard
constraints, $\beta$ a budget bound, and $\rho$ a verification requirement---for example,
``two independent sources must corroborate all contact details.'' The agent submits $I$ once. The
server returns a plan $P=(s_1,\dots,s_k)$ and a signature
$\sigma=\mathrm{Sign}_{sk}(H(P)\,\|\,H(I))$. Each step $s_j$ includes the provenance and content hash
of the artifact it references.

\begin{definition}[Acceptable plan]
The agent accepts $(P,\sigma)$ if and only if (i) $\mathrm{Verify}_{pk}(\sigma)$ succeeds, (ii) every
$s_j$ satisfies $C$ and the total cost does not exceed $\beta$, (iii) the attached provenance meets
$\rho$, and (iv) each artifact's hash still matches when the artifact is used. The final condition
closes the time-of-check/time-of-use window.
\end{definition}

\begin{proposition}[Round trips]
A stepwise interaction requires $\Theta(n)$ round trips for a chain of $n$ calls. Whole-intent
submission obtains the plan in $O(1)$ round trips and performs $O(k)$ local hash checks. It removes
the prediction step $\pi(t'\mid t,b)$ from the critical path.
\end{proposition}

\subsection{Limits}
Submitting $I$ reveals it to the server. A privacy-conscious design would therefore coarsen $g$ and
$C$, send $\beta$ as a bound rather than an exact budget, and prohibit uses of $I$ unrelated to the
request. Any advertising auction based on intent categories would require prominent disclosure.
Signatures establish who produced a statement, not whether the statement is true, so independent
corroboration under $\rho$ still carries the main trust burden.

\subsection{Possible implementation}
The current server accepts one tool call at a time through \texttt{/mcp}. A future version could add
a \texttt{submit\_intent} tool to \texttt{app/mcp.py} that accepts $I$ and returns $(P,\sigma)$; replace
the per-step predictor in \texttt{app/pipeline.py} with a \texttt{planner} module; and store a signing
key and hash-pinned artifacts alongside the rows in the \texttt{artifacts} table. Jev could still
classify intent categories or risk levels quickly, while a separate model composes the plan.

\section{Evidence about Jev and the broader model class}
\subsection{Available evidence}
Every available figure comes from the vendor or an individual developer report; \emph{none has been
independently verified}.
\begin{itemize}
\item TypeSafe reports response times of 70--500\,ms and input pricing of \$0.042 per million tokens.
      Its accuracy reference uses the average answer from two frontier models, and the start-up ran
      the evaluations itself \cite{heise}.
\item Other sources report ``20--200$\times$ faster'' and ``40--400$\times$ cheaper'' without naming
      the baselines \cite{latentspace}. Developers report gains of 5--18$\times$ on safety
      classification and 10--20$\times$ over Gemini on email classification. In the latter case,
      Gemini was slightly more accurate, while Jev also returned confidence scores \cite{techcrunch}.
\item The vendor announcement appears in \cite{typesafe}. Because Jev cannot emit free-form text,
      comparisons with general-purpose LLMs apply only to closed-output classification and routing.
\end{itemize}

\subsection{Generalization to other fast models}
Let model $M$ have latency $\lambda_M$, task accuracy $a_M$, and cost $c_M$. Fix an acceptability
threshold $a^\ast$, above which the prediction is useful, and let $\delta$ be the agent's latency
budget. Define the \emph{inline-feasible set}
\[
\mathcal F=\{M : a_M\ge a^\ast \;\wedge\; \lambda_M\le \varepsilon\,\delta\},\qquad \varepsilon\ll 1 .
\]
\begin{proposition}[Generality]
If $M\in\mathcal F$, prediction adds at most $\varepsilon\delta$ to the critical path, making the
prefetch/on-demand race available to any operator using $M$. If $M'$ is $k$ times faster than $M$ at
the same accuracy $a\ge a^\ast$, the operator gains
$\delta(\varepsilon-\varepsilon/k)$ of additional time. It can spend that time on more candidate
steps---increasing $|P_\tau|$ by lowering $\tau$---or on generation. Each order-of-magnitude speed
increase at acceptable quality therefore expands what the operator can prepare within $\delta$.
\end{proposition}
The argument does not depend on Jev's architecture. Small autoregressive models, distilled
classifiers, and non-autoregressive systems all qualify when $a_M\ge a^\ast$ and $\lambda_M$ is small
enough. Jev is one example of movement along the speed--quality frontier; the behavior described here
follows from that frontier, not from a particular vendor. Because the supply of sufficiently fast
predictors cannot be controlled, defenses must operate on the receiving side through provenance,
verification, and disclosure.

\section{Threat model for agent architectures}
The preceding sections describe a business mechanism. From a security perspective, they also
describe a class of vulnerabilities in agent architectures. An agent is exposed when it treats
output from an unauthenticated, attacker-reachable server as trustworthy input to its decisions.
Predictive prefetch can make that exposure cheaper and more reliable to exploit.

\subsection{Assets, actors, and trust boundary}
The \emph{assets} are the user's intent and data, the funds and credentials available to the agent,
and the agent's reputation with its user. The \emph{actors} are the end user (the principal), the
agent, an untrusted tool server $S$, and adversary $\mathcal A$, who operates $S$ or one of its pages.
The agent's context window is the \emph{trust boundary}. Everything that crosses this boundary from
$S$ is attacker-controlled data, even though an agent may use it to guide plans and actions.

\subsection{Vulnerability classes}
\begin{center}
\begin{tabular}{p{4.2cm}p{7.6cm}p{3cm}}
\hline
Class & How the mechanism contributes & Reference\\ \hline
Indirect prompt injection & Prefetched text arrives at the exact step when the agent reads it, making
  injected instructions more reliable. & OWASP LLM01\\
Excessive agency / confused deputy & The agent holds the user's authority, while tool content steers it
  to pay, message, or book on the attacker's behalf. & OWASP LLM06; CWE-441\\
Insufficient authenticity checks & The agent accepts new domains and lookalike sources because
  verification adds latency. & CWE-345\\
Time of check to time of use & An operator can replace a cached artifact between verification and use.
  & CWE-367\\
Predictable interaction as an oracle & Public call sequences make the next read predictable, allowing
  an adversary to prepare content instead of reacting after the fact. & Proposed in this paper\\
Supply-chain / namespace squatting & Cheap, automated domain claiming enables provider impersonation at
  scale. & CWE-1357 (reliance on an uncontrolled component)\\ \hline
\end{tabular}
\end{center}

\subsection{Why the race increases the risk}
Exploitation does not require a compromised host; it requires only that the agent \emph{read} the
attacker's content. The attacker pays for generation and a domain, while the defender pays the latency
cost of verification on every read. This asymmetry favors the attacker. The read-path race also
rewards operators that optimize for acceptance rather than truth.

\subsection{Qualitative severity}
We do not assign a CVSS score because the model does not describe a vulnerability in one specific
product. Qualitatively, the source is network-reachable, may require no authentication, and needs no
user interaction beyond ordinary tool use. The impact is bounded by the agent's authority and may be
high for agents that can send messages or make payments.

\subsection{Defender guidance}
\begin{enumerate}
\item Treat tool output as untrusted data rather than instructions, and separate it structurally from
      trusted instructions in the prompt.
\item Require provenance before using money, contact details, or credentials. Do not act on a single
      source seen for the first time.
\item Apply reputation and age gates to sources, and flag recently registered or first-seen domains.
\item Hash-pin artifacts between verification and use, and fetch them again if the hash changes.
\item Require human confirmation before irreversible or financial actions, regardless of model confidence.
\item Log and alert on new-domain sources, imperatives directed at the agent, content that changes
      between reads, and responses that arrive unusually quickly for the source.
\end{enumerate}

\section{Open questions}
Three questions remain. First, implementing a Jev-backed $\pi$ requires the actual API schema.
Second, measuring $H(\tau)$ requires representative agent traffic. Third, outreach to the Muse
(Meta) team has not yet been planned.

\begin{thebibliography}{9}
\bibitem{heise} heise online, ``AI model Jev to make machines decide faster,'' Sept.\ 2026.
  \url{https://www.heise.de/en/news/AI-model-Jev-to-make-machines-decide-faster-11457071.html}
\bibitem{latentspace} Latent Space, ``[AINews] Jev: a System One Model\dots,'' Sept.\ 2026.
  \url{https://www.latent.space/p/ainews-jev-a-system-one-model-that}
\bibitem{techcrunch} TechCrunch, ``A new kind of AI model from a ChatGPT inventor is thrilling developers,'' 18 Sept.\ 2026.
  \url{https://techcrunch.com/2026/09/18/a-new-kind-of-ai-model-from-a-chatgpt-inventor-is-thrilling-developers/}
\bibitem{typesafe} TypeSafe AI, ``Introducing System One Models \& Jev.''
  \url{https://typesafe.ai/blog/introducing-system-one-models-and-jev}
\end{thebibliography}

\end{document}
